\usepackage{kotex}
% 한글 폰트 — Apple SD Gothic Neo (macOS 표준, 받침/jongseong 안정)
\setmainhangulfont{Apple SD Gothic Neo}[BoldFont={* Bold}]
\setsanshangulfont{Apple SD Gothic Neo}[BoldFont={* Bold}]
\setmonohangulfont{Apple SD Gothic Neo}
% 영문 폰트 — 시스템 기본 사용
% underscore 자유 사용 (변수명 깨짐 방지)
\usepackage[strings]{underscore}
% 본문 줄간격 확대 — 가독성
\usepackage{setspace}
\onehalfspacing
% 단락 사이 여백 + 들여쓰기 0
\setlength{\parskip}{6pt}
\setlength{\parindent}{0pt}
% 목차 깊이 — N.M 까지만 (사용자 요청: 본문은 더 깊어도 됨)
\setcounter{tocdepth}{2}
\setcounter{secnumdepth}{4}
\usepackage{booktabs}
\usepackage{longtable}
\usepackage{array}
\usepackage{tabularx}
\usepackage{colortbl}
\usepackage{xcolor}
\usepackage{adjustbox}
% 모든 표에 셀 테두리 강제 (수정사항 #9 — 표 테두리)
\setlength{\arrayrulewidth}{0.4pt}
\renewcommand{\arraystretch}{1.18}
% Figure/Table 캡션 스타일
\usepackage{caption}
\captionsetup{font=small,labelfont=bf,labelsep=period}
% 표 폭 자동 축소 — Overfull hbox 방지
\usepackage{ragged2e}
% 표 전체 폰트 작게 + 셀 padding 축소
\AtBeginEnvironment{longtable}{\scriptsize\setlength{\tabcolsep}{3pt}}
\AtBeginEnvironment{tabular}{\scriptsize\setlength{\tabcolsep}{3pt}}
% longtable 의 모든 줄을 자동 ragged-right 처리 — 셀 내용이 길어도 줄바꿈
\setlength{\LTpre}{6pt}
\setlength{\LTpost}{6pt}
% Default column type 강화 — 모든 좌측 정렬 컬럼이 자동 줄바꿈하도록
% (pandoc longtable 의 'l' 컬럼은 wrap 안 됨, 'p{...}' 필요)
% → 컴파일 후 manual override 불가능하나, scriptsize 로 대부분 page width 안에 들어감
% 이미지 자동 크기 조정 — figure 페이지당 2-3개 배치되도록 폭 축소
\usepackage{graphicx}
\setkeys{Gin}{width=0.72\linewidth,keepaspectratio}
% figure float 옵션 — 본문 인접 배치 우선, page float 회피
\usepackage{float}
\makeatletter
\renewcommand{\fps@figure}{!htbp}
\makeatother
% 수학 기호 unicode → LaTeX 명령 매핑 (pandoc 이 math mode 를 unicode 로 재변환해도 안전)
\usepackage{newunicodechar}
\newunicodechar{≈}{\ensuremath{\approx}}
\newunicodechar{≥}{\ensuremath{\geq}}
\newunicodechar{≤}{\ensuremath{\leq}}
\newunicodechar{≠}{\ensuremath{\neq}}
\newunicodechar{ρ}{\ensuremath{\rho}}
\newunicodechar{ε}{\ensuremath{\varepsilon}}
\newunicodechar{μ}{\ensuremath{\mu}}
\newunicodechar{σ}{\ensuremath{\sigma}}
\newunicodechar{β}{\ensuremath{\beta}}
\newunicodechar{α}{\ensuremath{\alpha}}
\newunicodechar{γ}{\ensuremath{\gamma}}
\newunicodechar{Δ}{\ensuremath{\Delta}}
\newunicodechar{∈}{\ensuremath{\in}}
\newunicodechar{∉}{\ensuremath{\notin}}
\newunicodechar{∞}{\ensuremath{\infty}}
\newunicodechar{±}{\ensuremath{\pm}}
\newunicodechar{→}{\ensuremath{\rightarrow}}
\newunicodechar{←}{\ensuremath{\leftarrow}}
\newunicodechar{↑}{\ensuremath{\uparrow}}
\newunicodechar{↓}{\ensuremath{\downarrow}}
\newunicodechar{·}{\ensuremath{\cdot}}
\newunicodechar{×}{\ensuremath{\times}}
\newunicodechar{÷}{\ensuremath{\div}}
\newunicodechar{²}{\ensuremath{{}^2}}
\newunicodechar{³}{\ensuremath{{}^3}}
\newunicodechar{°}{\ensuremath{^\circ}}
